% Môn: Vật lý
% Lớp: 12
% Chương: Chương I. Vật lí nhiệt
% Bài: L12C1 Bài 6. Nhiệt hóa hơi riêng
% Số câu: 163
% App đọc file này trực tiếp — sửa rồi Commit trên GitHub, bấm Đồng bộ GitHub .tex

% L12C1 Bài 6. Nhiệt hóa hơi riêng





% ===== Câu 41 =====
% ID: L12C1B6-11-DS
% Mức: NB
\dangbt{Giải thích hiện tượng và ứng dụng liên quan đến sự hóa hơi}
\begin{ex}
% ID: L12C1B6-11-DS
% Mức: NB
Xét máy lạnh.
\choiceTF
{\True Môi chất bay hơi ở dàn lạnh và thu nhiệt.}
{\True Môi chất ngưng tụ ở dàn nóng và tỏa nhiệt.}
{\True Máy lạnh chuyển nhiệt từ nơi lạnh sang nơi nóng nhờ công của máy nén.}
{Máy lạnh tạo ra năng lượng từ hư không.}
\loigiai{
\begin{itemchoice}
\item (Đúng) Môi chất bay hơi ở dàn lạnh và thu nhiệt.
\item (Đúng) Môi chất ngưng tụ ở dàn nóng và tỏa nhiệt.
\item (Đúng) Máy lạnh chuyển nhiệt từ nơi lạnh sang nơi nóng nhờ công của máy nén.
\item (Sai) Máy lạnh tạo ra năng lượng từ hư không.
\end{itemchoice}
}
\end{ex}

% ===== Câu 44 =====
% ID: L12C1B6-11-TN
% Mức: VD
\dangbt{Bài toán tính toán nhiệt lượng cần thiết để hóa hơi hoàn toàn một lượng chất lỏng ở nhiệt độ sôi (áp dụng công thức Q =}
\begin{ex}
% ID: L12C1B6-11-TN
% Mức: VD
Biết nhiệt độ sôi, nhiệt dung riêng và nhiệt hóa hơi riêng của nước lần lượt là $100^\circ C$, $4200\ J/kg.K$ và $2,3.10^6\ J/kg$. Nhiệt lượng cần cung cấp để làm hóa hơi hoàn toàn $2\ kg$ nước từ nhiệt độ $20^\circ C$ là
\choice
{$2,6.10^6\ J$}
{\True $5,3.10^6\ J$}
{$26.10^6\ J$}
{$53.10^6\ J$}
\loigiai{
Nhiệt lượng cần cung cấp để làm hóa hơi hoàn toàn $2\,\text{kg}$ nước từ nhiệt độ $20^{\circ} {C}$ là:
${Q}={mc} \Delta {t}+{mL}=2 . 4200 .(100-20)+2 . 2,3 . 10^{6}=5272000 {~J} \approx 5,3 . 10^{6} {~J}$.
}
\end{ex}

% ===== Câu 45 =====
% ID: L12C1B6-12-DS
% Mức: TH
\dangbt{Giải thích hiện tượng và ứng dụng liên quan đến sự hóa hơi}
\begin{ex}
% ID: L12C1B6-12-DS
% Mức: TH
Xét nồi áp suất.
\choiceTF
{\True Áp suất trong nồi lớn hơn khí quyển.}
{\True Nhiệt độ sôi của nước tăng.}
{\True Thức ăn có thể chín nhanh hơn do nhiệt độ nấu cao.}
{Mở nắp ngay khi còn áp suất cao là an toàn.}
\loigiai{
\begin{itemchoice}
\item (Đúng) Áp suất trong nồi lớn hơn khí quyển.
\item (Đúng) Nhiệt độ sôi của nước tăng.
\item (Đúng) Thức ăn có thể chín nhanh hơn do nhiệt độ nấu cao.
\item (Sai) Mở nắp ngay khi còn áp suất cao là an toàn.
\end{itemchoice}
}
\end{ex}

% ===== Câu 49 =====
% ID: L12C1B6-13-DS
% Mức: VD
\begin{ex}
% ID: L12C1B6-13-DS
% Mức: VD
Xét hiện tượng bay hơi trong đời sống.
\choiceTF
{\True Bay hơi xảy ra ở mặt thoáng.}
{\True Nhiệt độ cao thường làm bay hơi nhanh hơn.}
{\True Gió có thể làm quần áo khô nhanh hơn.}
{Bay hơi luôn chỉ xảy ra ở 100 °C.}
\loigiai{
\begin{itemchoice}
\item (Đúng) Bay hơi xảy ra ở mặt thoáng.
\item (Đúng) Nhiệt độ cao thường làm bay hơi nhanh hơn. (Vì): ay hơi có thể xảy ra ở mọi nhiệt độ; ba yếu tố thường gặp là nhiệt độ, diện tích mặt thoáng và gió
\item (Đúng) Gió có thể làm quần áo khô nhanh hơn.
\item (Sai) Bay hơi luôn chỉ xảy ra ở 100 °C.
\end{itemchoice}
}
\end{ex}

% ===== Câu 50 =====
% ID: L12C1B6-13-TL
% Mức: NB
\begin{bt}
% ID: L12C1B6-13-TL
% Mức: NB
Giải thích vì sao thức ăn chín nhanh trong nồi áp suất và nêu lưu ý an toàn.
\loigiai{
Áp suất cao làm nhiệt độ sôi của nước cao hơn 100 ° $C$, nên thức ăn được nấu ở nhiệt độ cao và chín nhanh. Phải kiểm tra van, không mở nắp khi còn áp suất cao và xả áp đúng hướng dẫn.
}
\end{bt}

% ===== Câu 51 =====
% ID: L12C1B6-13-TLN
% Mức: NB
\begin{ex}
% ID: L12C1B6-13-TLN
% Mức: NB
Cồn bay hơi $20\,\text{g}, L=0,9.10^6\,\text{J}$ /kg. Nhiệt lượng thu theo kJ là bao nhiêu? Chỉ ghi số.
\shortans{18}
\loigiai{
$m=0,020 kg; Q=0,9.10^6.0,020=18 kJ.$
}
\end{ex}

% ===== Câu 53 =====
% ID: L12C1B6-14-DS
% Mức: VD
\begin{ex}
% ID: L12C1B6-14-DS
% Mức: VD
Các biện pháp làm khô.
\choiceTF
{\True Phơi rộng quần áo làm tăng diện tích bay hơi.}
{\True Thông gió mang hơi ẩm đi.}
{Không khí ẩm bão hòa làm bay hơi nhanh hơn không khí khô.}
{\True Sấy ấm thường tăng tốc độ bay hơi.}
\loigiai{
\begin{itemchoice}
\item (Đúng) Phơi rộng quần áo làm tăng diện tích bay hơi.
\item (Đúng) Thông gió mang hơi ẩm đi.
\item (Sai) Không khí ẩm bão hòa làm bay hơi nhanh hơn không khí khô.
\item (Đúng) Sấy ấm thường tăng tốc độ bay hơi.
\end{itemchoice}
}
\end{ex}

% ===== Câu 54 =====
% ID: L12C1B6-14-TL
% Mức: TH
\begin{bt}
% ID: L12C1B6-14-TL
% Mức: TH
Giải thích vì sao cồn sát khuẩn trên da gây cảm giác mát và vì sao quạt làm cảm giác mát rõ hơn.
\loigiai{
Cồn bay hơi cần năng lượng nên lấy nhiệt từ da, làm da giảm nhiệt độ. Quạt mang lớp không khí giàu hơi cồn đi, tăng tốc độ bay hơi nên tốc độ lấy nhiệt tăng và cảm giác mát rõ hơn.
}
\end{bt}

% ===== Câu 55 =====
% ID: L12C1B6-14-TLN
% Mức: NB
\begin{ex}Một thiết bị lấy khỏi phòng $900 \text{ kJ}$ nhờ $0,50 \text{ kg}$ môi chất bay hơi. Tính nhiệt hóa hơi riêng $L$ của môi chất theo đơn vị $10^5 \text{ J/kg}$, chỉ ghi số.
\shortans{18}
\loigiai{
Nhiệt lượng mà thiết bị lấy khỏi phòng là $Q = 900 \text{ kJ} = 900 \times 10^3 \text{ J}$.
Khối lượng môi chất bay hơi là $m = 0,50 \text{ kg}$.
Nhiệt hóa hơi riêng $L$ được tính bằng công thức $L = \frac{Q}{m}$.
Thay số vào, ta có:
$L = \frac{900 \times 10^3 \text{ J}}{0,50 \text{ kg}} = 1800 \times 10^3 \text{ J/kg} = 1,8 \times 10^6 \text{ J/kg}$.
Để biểu diễn $L$ theo đơn vị $10^5 \text{ J/kg}$, ta chuyển đổi:
$L = 1,8 \times 10^6 \text{ J/kg} = 18 \times 10^5 \text{ J/kg}$.
Vậy, giá trị của $L$ theo đơn vị $10^5 \text{ J/kg}$ là $18$.
}
\end{ex}

% ===== Câu 57 =====
% ID: L12C1B6-15-DS
% Mức: VDC
\begin{ex}
% ID: L12C1B6-15-DS
% Mức: VDC
Về bỏng hơi nước.
\choiceTF
{\True Hơi 100 ° $C$ có thể ngưng tụ trên da.}
{\True Ngưng tụ tỏa nhiệt Lm.}
{\True Cùng khối lượng, hơi 100 ° $C$ thường truyền nhiều nhiệt hơn nước 100 ° $C$ khi cùng hạ đến nhiệt độ da.}
{Hơi nước không gây bỏng vì là chất khí.}
\loigiai{
\begin{itemchoice}
\item (Đúng) Hơi 100 ° $C$ có thể ngưng tụ trên da.
\item (Đúng) Ngưng tụ tỏa nhiệt Lm.
\item (Đúng) Cùng khối lượng, hơi 100 ° $C$ thường truyền nhiều nhiệt hơn nước 100 ° $C$ khi cùng hạ đến nhiệt độ da.
\item (Sai) Hơi nước không gây bỏng vì là chất khí.
\end{itemchoice}
}
\end{ex}

% ===== Câu 58 =====
% ID: L12C1B6-15-TL
% Mức: TH
\begin{bt}
% ID: L12C1B6-15-TL
% Mức: TH
Phân tích ảnh hưởng của nhiệt độ, diện tích mặt thoáng, gió và độ ẩm đến tốc độ bay hơi.
\loigiai{
Nhiệt độ cao làm phân tử dễ thoát khỏi mặt thoáng; diện tích lớn tăng số phân tử có thể thoát; gió mang hơi đi; độ ẩm thấp giúp không khí nhận thêm hơi. Vì vậy tăng ba yếu tố đầu và giảm độ ẩm thường làm bay hơi nhanh.
}
\end{bt}

% ===== Câu 59 =====
% ID: L12C1B6-15-TLN
% Mức: TH
\begin{ex}
% ID: L12C1B6-15-TLN
% Mức: TH
Nước bay hơi $50\,\text{g}$ khỏi sân, L=2,3.10^ $6\,\text{J}$ /kg. Nhiệt lượng lấy từ môi trường theo kJ là bao nhiêu? Chỉ ghi số.
\shortans{115}
\loigiai{
$Q=2,3.10^6.0,050=115 kJ.$
}
\end{ex}

% ===== Câu 62 =====
% ID: L12C1B6-16-TL
% Mức: VD
\begin{bt}
% ID: L12C1B6-16-TL
% Mức: VD
Giải thích vì sao bỏng hơi nước 100 ° $C$ nguy hiểm hơn bỏng nước 100 ° $C$ với cùng khối lượng.
\loigiai{
Hơi nước khi gặp da ngưng tụ và tỏa thêm nhiệt lượng Lm, sau đó nước ngưng còn nguội từ 100 ° $C$ xuống nhiệt độ da. Nước 100 ° $C$ chỉ có phần nguội đi, nên hơi truyền năng lượng lớn hơn.
}
\end{bt}

% ===== Câu 63 =====
% ID: L12C1B6-16-TLN
% Mức: VD
\begin{ex}
% ID: L12C1B6-16-TLN
% Mức: VD
Một dàn lạnh cho $0,20\,\text{kg}$ môi chất bay hơi mỗi phút, L=1,5.10^ $5\,\text{J}$ /kg. Nhiệt thu mỗi phút theo kJ là bao nhiêu? Chỉ ghi số.
\shortans{30}
\loigiai{
$Q=Lm=1,5.10^5.0,20=30000 J=30 kJ.$
}
\end{ex}

% ===== Câu 66 =====
% ID: L12C1B6-17-TL
% Mức: VD
\begin{bt}
% ID: L12C1B6-17-TL
% Mức: VD
Một người nói “phơi trong phòng kín thật nóng luôn nhanh khô nhất”. Hãy nhận xét.
\loigiai{
Không hoàn toàn đúng. Nhiệt độ cao thúc đẩy bay hơi, nhưng phòng kín nhanh trở nên ẩm, làm tốc độ bay hơi giảm. Cần thêm thông gió hoặc hút ẩm để mang hơi nước đi.
}
\end{bt}

% ===== Câu 67 =====
% ID: L12C1B6-17-TLN
% Mức: VDC
\begin{ex}
% ID: L12C1B6-17-TLN
% Mức: VDC
Hơi nước $0,10\,\text{kg}$ ở 100 ° $C$ ngưng tụ, lấy L=2,3.10^ $6\,\text{J}$ /kg. Riêng giai đoạn ngưng tụ tỏa bao nhiêu kJ? Chỉ ghi số.
\shortans{230}
\loigiai{
$Q=Lm=2,3.10^6.0,10=230 kJ.$
}
\end{ex}

% ===== Câu 70 =====
% ID: L12C1B6-18-TL
% Mức: VDC
\begin{bt}
% ID: L12C1B6-18-TL
% Mức: VDC
Trình bày vai trò của hóa hơi và ngưng tụ trong chu trình làm lạnh của điều hòa.
\loigiai{
Ở dàn lạnh, môi chất áp suất thấp bay hơi và hấp thụ nhiệt trong phòng. Máy nén thực hiện công; ở dàn nóng, môi chất ngưng tụ và tỏa nhiệt ra ngoài. Van tiết lưu hạ áp để chu trình lặp lại.
}
\end{bt}

% ===== Câu 71 =====
% ID: L12C1B6-18-TLN
% Mức: VDC
\begin{ex}
% ID: L12C1B6-18-TLN
% Mức: VDC
Hơi nước ngưng tụ tỏa 460 kJ, L=2,3.10^ $6\,\text{J}$ /kg. Khối lượng hơi theo gam là bao nhiêu? Chỉ ghi số.
\shortans{200}
\loigiai{
$m=460000/(2,3.10^6)=0,20 kg=200 g.$
}
\end{ex}

% ===== Câu 88 =====
% ID: L12C1B6-22-TN
% Mức: NB
\begin{ex}
% ID: L12C1B6-22-TN
% Mức: NB
Cồn bôi lên da tạo cảm giác mát chủ yếu vì
\choice
{cồn đông đặc}
{\True cồn bay hơi thu nhiệt từ da}
{cồn ngưng tụ tỏa nhiệt}
{nhiệt độ da tự tăng}
\loigiai{
Bay hơi là quá trình thu nhiệt; cồn lấy nhiệt từ da nên da mát.
}
\end{ex}

% ===== Câu 92 =====
% ID: L12C1B6-23-TN
% Mức: NB
\begin{ex}
% ID: L12C1B6-23-TN
% Mức: NB
Trong tủ lạnh, môi chất lạnh làm lạnh khoang khi
\choice
{\True bay hơi và thu nhiệt}
{ngưng tụ và thu nhiệt}
{đông đặc và thu nhiệt}
{tăng áp không trao đổi nhiệt}
\loigiai{
Môi chất bay hơi ở dàn lạnh, hấp thụ nhiệt từ khoang.
}
\end{ex}

% ===== Câu 96 =====
% ID: L12C1B6-24-TN
% Mức: NB
\begin{ex}
% ID: L12C1B6-24-TN
% Mức: NB
Hơi nước ngưng tụ trên thành cốc lạnh làm thành cốc
\choice
{\True nhận nhiệt từ hơi}
{không trao đổi nhiệt}
{chỉ giảm khối lượng}
{hóa hơi nhanh hơn}
\loigiai{
Khi ngưng tụ, hơi nước tỏa nhiệt cho thành cốc và môi trường.
}
\end{ex}

% ===== Câu 100 =====
% ID: L12C1B6-25-TN
% Mức: TH
\begin{ex}
% ID: L12C1B6-25-TN
% Mức: TH
Quần áo nhanh khô hơn khi có gió vì gió
\choice
{làm giảm diện tích mặt thoáng}
{\True mang lớp hơi ẩm sát bề mặt đi}
{làm nước ngưng tụ}
{làm $L$ của nước bằng 0}
\loigiai{
Gió làm giảm độ ẩm cục bộ gần bề mặt, tăng tốc độ bay hơi.
}
\end{ex}

% ===== Câu 101 =====
% ID: L12C1B6-26-DS
% Mức: NB
\dangbt{Tính nhiệt lượng, khối lượng và nhiệt hóa hơi riêng bằng Q = Lm}
\begin{ex}
% ID: L12C1B6-26-DS
% Mức: NB
Để xác định $L$ bằng phép đo $Q$ và m.
\choiceTF
{\True Có thể dùng L=Q/m.}
{\True $Q$ phải là phần năng lượng hữu ích cho hóa hơi.}
{\True m cần đổi về kg nếu tính $L$ theo J/kg.}
{Sai số của $Q$ và m không ảnh hưởng kết quả.}
\loigiai{
\begin{itemchoice}
\item (Đúng) Có thể dùng L=Q/m.
\item (Đúng) $Q$ phải là phần năng lượng hữu ích cho hóa hơi. (Vì): a phát biểu đầu là điều kiện tính; đại lượng đo có sai số nên $L$ cũng có sai số
\item (Đúng) m cần đổi về kg nếu tính $L$ theo J/kg.
\item (Sai) Sai số của $Q$ và m không ảnh hưởng kết quả.
\end{itemchoice}
}
\end{ex}

% ===== Câu 104 =====
% ID: L12C1B6-26-TN
% Mức: TH
\dangbt{Giải thích hiện tượng và ứng dụng liên quan đến sự hóa hơi}
\begin{ex}
% ID: L12C1B6-26-TN
% Mức: TH
Trải nước lên sân nóng giúp sân mát đi vì nước
\choice
{\True bay hơi thu nhiệt}
{ngưng tụ tỏa nhiệt}
{làm tăng bức xạ Mặt Trời}
{có nhiệt độ sôi bằng 0 ° $C$}
\loigiai{
Nước bay hơi lấy nhiệt từ sân và không khí gần sân.
}
\end{ex}

% ===== Câu 105 =====
% ID: L12C1B6-27-DS
% Mức: TH
\dangbt{Tính nhiệt lượng, khối lượng và nhiệt hóa hơi riêng bằng Q = Lm}
\begin{ex}
% ID: L12C1B6-27-DS
% Mức: TH
Hai mẫu cùng chất ở nhiệt độ sôi có m_A=2m_B.
\choiceTF
{\True $Q_A=2Q_B.$}
{\True $Q_A/Q_B=2.$}
{$L_A=2L_B.$}
{\True Nếu ngưng tụ thì độ lớn nhiệt lượng của $A$ cũng gấp đôi B.}
\loigiai{
\begin{itemchoice}
\item (Đúng) $Q_A=2Q_B.$.
\item (Đúng) $Q_A/Q_B=2.$.
\item (Sai) $L_A=2L_B.$. (Vì): ùng chất nên L_A=L_ $B$; Q=Lm
\item (Đúng) Nếu ngưng tụ thì độ lớn nhiệt lượng của $A$ cũng gấp đôi B.
\end{itemchoice}
}
\end{ex}

% ===== Câu 108 =====
% ID: L12C1B6-27-TN
% Mức: TH
\dangbt{Giải thích hiện tượng và ứng dụng liên quan đến sự hóa hơi}
\begin{ex}
% ID: L12C1B6-27-TN
% Mức: TH
Trong điều hòa không khí, bộ phận đặt trong phòng cần tạo điều kiện để môi chất
\choice
{\True bay hơi}
{ngưng tụ}
{đông đặc}
{nóng chảy}
\loigiai{
Tại dàn lạnh, môi chất bay hơi và hấp thụ nhiệt trong phòng.
}
\end{ex}

% ===== Câu 109 =====
% ID: L12C1B6-28-DS
% Mức: VD
\dangbt{Tính nhiệt lượng, khối lượng và nhiệt hóa hơi riêng bằng Q = Lm}
\begin{ex}
% ID: L12C1B6-28-DS
% Mức: VD
Hóa hơi $0,50\,\text{kg}$ nước ở 100 ° $C$, lấy L=2,3.10^ $6\,\text{J}$ /kg.
\choiceTF
{\True $Q=1,15.10^6 J.$}
{\True $Q=1150 kJ.$}
{\True Nếu chỉ hóa hơi $0,25\,\text{kg}$ thì $Q$ giảm một nửa.}
{$Q$ không phụ thuộc khối lượng.}
\loigiai{
\begin{itemchoice}
\item (Đúng) $Q=1,15.10^6 J.$.
\item (Đúng) $Q=1150 kJ.$.
\item (Đúng) Nếu chỉ hóa hơi $0,25\,\text{kg}$ thì $Q$ giảm một nửa.
\item (Sai) $Q$ không phụ thuộc khối lượng.
\end{itemchoice}
}
\end{ex}

% ===== Câu 112 =====
% ID: L12C1B6-28-TN
% Mức: VD
\dangbt{Giải thích hiện tượng và ứng dụng liên quan đến sự hóa hơi}
\begin{ex}
% ID: L12C1B6-28-TN
% Mức: VD
Nồi áp suất làm thức ăn chín nhanh vì
\choice
{áp suất giảm, nhiệt độ sôi giảm}
{\True áp suất tăng, nhiệt độ sôi tăng}
{nước không cần thu nhiệt}
{nhiệt hóa hơi bằng 0}
\loigiai{
Áp suất cao làm nhiệt độ sôi của nước tăng, thức ăn được nấu ở nhiệt độ cao hơn.
}
\end{ex}

% ===== Câu 113 =====
% ID: L12C1B6-29-DS
% Mức: VD
\dangbt{Tính nhiệt lượng, khối lượng và nhiệt hóa hơi riêng bằng Q = Lm}
\begin{ex}
% ID: L12C1B6-29-DS
% Mức: VD
Kiểm tra đơn vị trong Q=Lm.
\choiceTF
{\True J/kg nhân kg cho J.}
{\True kJ có thể đổi sang $J$ bằng cách nhân 1000.}
{Gam có thể thay trực tiếp cho kilôgam mà không đổi.}
{\True Nếu $Q$ tính kJ và m kg thì $L$ tính kJ/kg.}
\loigiai{
\begin{itemchoice}
\item (Đúng) J/kg nhân kg cho J.
\item (Đúng) kJ có thể đổi sang $J$ bằng cách nhân 1000.
\item (Sai) Gam có thể thay trực tiếp cho kilôgam mà không đổi.
\item (Đúng) Nếu $Q$ tính kJ và m kg thì $L$ tính kJ/kg.
\end{itemchoice}
}
\end{ex}

% ===== Câu 116 =====
% ID: L12C1B6-29-TN
% Mức: VD
\dangbt{Giải thích hiện tượng và ứng dụng liên quan đến sự hóa hơi}
\begin{ex}
% ID: L12C1B6-29-TN
% Mức: VD
Muốn tăng tốc độ bay hơi của nước, biện pháp phù hợp là
\choice
{giảm diện tích mặt thoáng}
{tăng độ ẩm không khí}
{\True tăng nhiệt độ và thông gió}
{đậy kín bình}
\loigiai{
Nhiệt độ cao, diện tích thoáng lớn và gió thúc đẩy bay hơi.
}
\end{ex}

% ===== Câu 117 =====
% ID: L12C1B6-30-DS
% Mức: VDC
\dangbt{Tính nhiệt lượng, khối lượng và nhiệt hóa hơi riêng bằng Q = Lm}
\begin{ex}
% ID: L12C1B6-30-DS
% Mức: VDC
Một chất có L=1,2.10^ $6\,\text{J}$ /kg; cung cấp 360 kJ chỉ cho hóa hơi.
\choiceTF
{\True Khối lượng hóa hơi là $0,30\,\text{kg}$.}
{\True Khối lượng hóa hơi là $300\,\text{g}$.}
{Nếu hiệu suất chỉ 50% thì cùng năng lượng nguồn vẫn hóa hơi $0,30\,\text{kg}$.}
{\True $L=Q/m.$}
\loigiai{
\begin{itemchoice}
\item (Đúng) Khối lượng hóa hơi là $0,30\,\text{kg}$.
\item (Đúng) Khối lượng hóa hơi là $300\,\text{g}$.
\item (Sai) Nếu hiệu suất chỉ 50% thì cùng năng lượng nguồn vẫn hóa hơi $0,30\,\text{kg}$.
\item (Đúng) $L=Q/m.$.
\end{itemchoice}
}
\end{ex}

% ===== Câu 120 =====
% ID: L12C1B6-30-TN
% Mức: VDC
\dangbt{Giải thích hiện tượng và ứng dụng liên quan đến sự hóa hơi}
\begin{ex}
% ID: L12C1B6-30-TN
% Mức: VDC
Bỏng do hơi nước 100 ° $C$ thường nặng hơn do nước 100 ° $C$ vì hơi nước
\choice
{có khối lượng bằng 0}
{\True khi ngưng tụ còn tỏa nhiệt Lm}
{không truyền nhiệt}
{làm da bay hơi tức thì}
\loigiai{
Hơi ngưng tụ trên da tỏa thêm nhiệt hóa hơi lớn.
}
\end{ex}

% ===== Câu 121 =====
% ID: L12C1B6-31-TL
% Mức: NB
\dangbt{Tính nhiệt lượng, khối lượng và nhiệt hóa hơi riêng bằng Q = Lm}
\begin{bt}
% ID: L12C1B6-31-TL
% Mức: NB
Hai chất $A$ và $B$ cùng khối lượng $0,30\,\text{kg}$ có L_A=2,0.10^ $6\,\text{J}$ /kg, L_B=0,9.10^ $6\,\text{J}$ /kg. Tính và so sánh nhiệt lượng hóa hơi.
\loigiai{
Q_A=2,0.10^6.0,30=600 kJ; Q_B=0,9.10^6.0,30=270 kJ. Chất $A$ cần nhiều hơn 330 kJ, tỉ số Q_A/Q_B=2,22.
}
\end{bt}

% ===== Câu 122 =====
% ID: L12C1B6-31-TLN
% Mức: NB
\begin{ex}
% ID: L12C1B6-31-TLN
% Mức: NB
Cung cấp 540 kJ để hóa hơi chất có L=1,8.10^ $6\,\text{J}$ /kg. Tính khối lượng theo gam, chỉ ghi số.
\shortans{300}
\loigiai{
$m=540000/(1,8.10^6)=0,30 kg=300 g.$
}
\end{ex}

% ===== Câu 123 =====
% ID: L12C1B6-31-TN
% Mức: VDC
\dangbt{Giải thích hiện tượng và ứng dụng liên quan đến sự hóa hơi}
\begin{ex}
% ID: L12C1B6-31-TN
% Mức: VDC
Sự bay hơi khác sự sôi ở chỗ bay hơi
\choice
{chỉ xảy ra ở nhiệt độ sôi}
{\True xảy ra trên mặt thoáng ở mọi nhiệt độ}
{xảy ra trong toàn khối ở nhiệt độ xác định}
{không thu nhiệt}
\loigiai{
Bay hơi xảy ra trên mặt thoáng và có thể ở mọi nhiệt độ.
}
\end{ex}

% ===== Câu 124 =====
% ID: L12C1B6-32-TL
% Mức: TH
\dangbt{Tính nhiệt lượng, khối lượng và nhiệt hóa hơi riêng bằng Q = Lm}
\begin{bt}
% ID: L12C1B6-32-TL
% Mức: TH
Tính nhiệt lượng cần để hóa hơi hoàn toàn $0,75\,\text{kg}$ nước ở 100 ° $C$, lấy L=2,3.10^ $6\,\text{J}$ /kg.
\loigiai{
Áp dụng Q=Lm=2,3.10^6.0,75=1,725.10^ $6\,\text{J}$ =1725 kJ.
}
\end{bt}

% ===== Câu 125 =====
% ID: L12C1B6-32-TLN
% Mức: NB
\begin{ex}
% ID: L12C1B6-32-TLN
% Mức: NB
Q_A=720 kJ dùng hóa hơi $0,40\,\text{kg}$ chất A. Tính L_ $A$ theo đơn vị 10^ $6\,\text{J}$ /kg, chỉ ghi số.
\shortans{01/08/2026}
\loigiai{
$L=720000/0,40=1,8.10^6 J/kg.$
}
\end{ex}

% ===== Câu 127 =====
% ID: L12C1B6-33-TL
% Mức: TH
\begin{bt}
% ID: L12C1B6-33-TL
% Mức: TH
Một mẫu hơi $0,50\,\text{kg}$ ngưng tụ hoàn toàn, L=1,6.10^ $6\,\text{J}$ /kg. Tính nhiệt lượng tỏa ra và nêu dấu của $Q$ nếu quy ước hệ nhận nhiệt dương.
\loigiai{
Độ lớn nhiệt lượng tỏa ra là Lm=1,6.10^6.0,50=8,0.10^ $5\,\text{J}$. Theo quy ước hệ nhận nhiệt dương, hơi tỏa nhiệt nên Q=-8,0.10^ $5\,\text{J}$.
}
\end{bt}

% ===== Câu 128 =====
% ID: L12C1B6-33-TLN
% Mức: TH
\begin{ex}
% ID: L12C1B6-33-TLN
% Mức: TH
Hóa hơi $250\,\text{g}$ chất lỏng cần 200 kJ. Tính $L$ theo đơn vị 10^ $5\,\text{J}$ /kg, chỉ ghi số.
\shortans{8}
\loigiai{
$m=0,25 kg; L=200000/0,25=8,0.10^5 J/kg.$
}
\end{ex}

% ===== Câu 130 =====
% ID: L12C1B6-34-TL
% Mức: VD
\begin{bt}
% ID: L12C1B6-34-TL
% Mức: VD
Một bình mất $0,24\,\text{kg}$ chất lỏng do sôi; năng lượng hữu ích đã truyền là 432 kJ. Xác định nhiệt hóa hơi riêng.
\loigiai{
$L=Q/m=432000/0,24=1,8.10^6 J/kg.$
}
\end{bt}

% ===== Câu 131 =====
% ID: L12C1B6-34-TLN
% Mức: VD
\begin{ex}
% ID: L12C1B6-34-TLN
% Mức: VD
Ngưng tụ $0,12\,\text{kg}$ hơi nước, lấy L=2,3.10^ $6\,\text{J}$ /kg. Tính nhiệt lượng tỏa ra theo kJ, chỉ ghi số.
\shortans{276}
\loigiai{
$Q=Lm=276000 J=276 kJ.$
}
\end{ex}

% ===== Câu 133 =====
% ID: L12C1B6-35-TL
% Mức: VD
\begin{bt}
% ID: L12C1B6-35-TL
% Mức: VD
Chứng minh cách xác định $L$ từ đồ thị $Q$ theo m của cùng một chất ở nhiệt độ sôi.
\loigiai{
Từ Q=Lm, đồ thị $Q$ theo m là đường thẳng qua gốc. Hệ số góc ΔQ/Δm chính là L. Đo nhiều cặp (m,Q), vẽ đường phù hợp và lấy độ dốc giúp giảm ảnh hưởng sai số ngẫu nhiên.
}
\end{bt}

% ===== Câu 134 =====
% ID: L12C1B6-35-TLN
% Mức: VDC
\begin{ex}
% ID: L12C1B6-35-TLN
% Mức: VDC
Hóa hơi $0,18\,\text{kg}$ nước ở nhiệt độ sôi với L=2,3.10^ $6\,\text{J}$ /kg. Tính $Q$ theo kJ, chỉ ghi số.
\shortans{414}
\loigiai{
$Q=2,3.10^6.0,18=414000 J=414 kJ.$
}
\end{ex}

% ===== Câu 136 =====
% ID: L12C1B6-36-TL
% Mức: VDC
\begin{bt}
% ID: L12C1B6-36-TL
% Mức: VDC
Cung cấp 920 kJ để hóa hơi nước ở 100 °C. Lấy L=2,3.10^ $6\,\text{J}$ /kg. Tính khối lượng nước hóa hơi.
\loigiai{
$m=Q/L=920000/(2,3.10^6)=0,40 kg.$
}
\end{bt}

% ===== Câu 137 =====
% ID: L12C1B6-36-TLN
% Mức: VDC
\begin{ex}
% ID: L12C1B6-36-TLN
% Mức: VDC
Một nguồn cung cấp 1,2 MJ hữu ích cho chất có L=2,0.10^ $6\,\text{J}$ /kg. Khối lượng hóa hơi theo gam là bao nhiêu? Chỉ ghi số.
\shortans{600}
\loigiai{
$m=1,2.10^6/(2,0.10^6)=0,60 kg=600 g.$
}
\end{ex}

% ===== Câu 154 =====
% ID: L12C1B6-52-TN
% Mức: NB
\begin{ex}
% ID: L12C1B6-52-TN
% Mức: NB
Hóa hơi hoàn toàn $0,2\,\text{kg}$ chất lỏng ở nhiệt độ sôi, biết L=2,3.10^ $6\,\text{J}$ /kg. Nhiệt lượng cần cung cấp là
\choice
{46 kJ}
{\True 460 kJ}
{4600 kJ}
{460 $J$}
\loigiai{
$Q=Lm=2,3.10^6.0,2=460 kJ.$
}
\end{ex}

% ===== Câu 155 =====
% ID: L12C1B6-53-TN
% Mức: NB
\begin{ex}
% ID: L12C1B6-53-TN
% Mức: NB
Hóa hơi hoàn toàn $0,25\,\text{kg}$ chất lỏng ở nhiệt độ sôi, biết L=2,4.10^ $6\,\text{J}$ /kg. Nhiệt lượng cần cung cấp là
\choice
{60 kJ}
{\True 600 kJ}
{6000 kJ}
{600 $J$}
\loigiai{
$Q=Lm=2,4.10^6.0,25=600 kJ.$
}
\end{ex}

% ===== Câu 156 =====
% ID: L12C1B6-54-TN
% Mức: NB
\begin{ex}
% ID: L12C1B6-54-TN
% Mức: NB
Hóa hơi $0,30\,\text{kg}$ chất lỏng cần 270 kJ. Nhiệt hóa hơi riêng là
\choice
{$9,0.10^4 J/kg$}
{\True $9,0.10^5 J/kg$}
{$9,0.10^6 J/kg$}
{$8,1.10^4 J/kg$}
\loigiai{
$L=Q/m=270000/0,30=9,0.10^5 J/kg.$
}
\end{ex}

% ===== Câu 157 =====
% ID: L12C1B6-55-TN
% Mức: TH
\begin{ex}
% ID: L12C1B6-55-TN
% Mức: TH
Hóa hơi hoàn toàn $0,35\,\text{kg}$ chất lỏng ở nhiệt độ sôi, biết L=2.10^ $6\,\text{J}$ /kg. Nhiệt lượng cần cung cấp là
\choice
{70 kJ}
{\True 700 kJ}
{7000 kJ}
{700 $J$}
\loigiai{
$Q=Lm=2.10^6.0,35=700 kJ.$
}
\end{ex}

% ===== Câu 158 =====
% ID: L12C1B6-56-TN
% Mức: TH
\begin{ex}
% ID: L12C1B6-56-TN
% Mức: TH
Cần 460 kJ để hóa hơi chất lỏng có L=2,3.10^ $6\,\text{J}$ /kg. Khối lượng chất lỏng là
\choice
{$0,02\,\text{kg}$}
{\True $0,2\,\text{kg}$}
{$2\,\text{kg}$}
{$200\,\text{kg}$}
\loigiai{
$m=Q/L=460000/(2,3.10^6)=0,2 kg.$
}
\end{ex}

% ===== Câu 159 =====
% ID: L12C1B6-57-TN
% Mức: TH
\begin{ex}
% ID: L12C1B6-57-TN
% Mức: TH
Hai chất $A$, $B$ có cùng khối lượng hóa hơi; L_A=2L_B. Tỉ số Q_A/Q_ $B$ là
\choice
{01/02/2026}
{1}
{\True 2}
{4}
\loigiai{
Q=Lm nên với cùng m, Q_A/Q_B=L_A/L_B=2.
}
\end{ex}

% ===== Câu 160 =====
% ID: L12C1B6-58-TN
% Mức: VD
\begin{ex}
% ID: L12C1B6-58-TN
% Mức: VD
Hóa hơi hoàn toàn $0,4\,\text{kg}$ chất lỏng ở nhiệt độ sôi, biết L=0,9.10^ $6\,\text{J}$ /kg. Nhiệt lượng cần cung cấp là
\choice
{36 kJ}
{\True 360 kJ}
{3600 kJ}
{360 $J$}
\loigiai{
$Q=Lm=0,9.10^6.0,4=360 kJ.$
}
\end{ex}

% ===== Câu 161 =====
% ID: L12C1B6-59-TN
% Mức: VD
\begin{ex}
% ID: L12C1B6-59-TN
% Mức: VD
Cần 600 kJ để hóa hơi chất lỏng có L=2.10^ $6\,\text{J}$ /kg. Khối lượng chất lỏng là
\choice
{$0,03\,\text{kg}$}
{\True $0,3\,\text{kg}$}
{$3\,\text{kg}$}
{$300\,\text{kg}$}
\loigiai{
$m=Q/L=600000/(2.10^6)=0,3 kg.$
}
\end{ex}

% ===== Câu 162 =====
% ID: L12C1B6-60-TN
% Mức: VDC
\begin{ex}
% ID: L12C1B6-60-TN
% Mức: VDC
Hóa hơi hoàn toàn $0,6\,\text{kg}$ chất lỏng ở nhiệt độ sôi, biết L=1,5.10^ $6\,\text{J}$ /kg. Nhiệt lượng cần cung cấp là
\choice
{90 kJ}
{\True 900 kJ}
{9000 kJ}
{900 $J$}
\loigiai{
$Q=Lm=1,5.10^6.0,6=900 kJ.$
}
\end{ex}

% ===== Câu 163 =====
% ID: L12C1B6-61-TN
% Mức: VDC
\begin{ex}
% ID: L12C1B6-61-TN
% Mức: VDC
Cần 450 kJ để hóa hơi chất lỏng có L=1,5.10^ $6\,\text{J}$ /kg. Khối lượng chất lỏng là
\choice
{$0,03\,\text{kg}$}
{\True $0,3\,\text{kg}$}
{$3\,\text{kg}$}
{$300\,\text{kg}$}
\loigiai{
$m=Q/L=450000/(1,5.10^6)=0,3 kg.$
}
\end{ex}
